\documentclass[../Main.tex]{subfiles}
\begin{document}
Chapter 2 presents a comprehensive survey and analysis of the system requirements. The chapter begins with a status survey to establish the project context, followed by a functional overview that illustrates the system's scope through general and detailed use case diagrams as well as business processes. Subsequently, specific features are examined in depth in the functional description section. Finally, the chapter concludes by outlining the essential non-functional requirements necessary to ensure the system's performance and quality standards.

\section{Status survey}
\label{section:2.1}

A thorough survey of the current landscape is essential to identify the precise software requirements and ensure the proposed solution effectively addresses practical challenges. This survey analyzes data from three primary sources: target users in flood-prone areas, existing emergency response systems, and similar location-sharing applications on the market.

% \subsection{User and Practical Context Survey}
% \label{subsection:2.1.1}

% A fundamental source of software requirements stems from analyzing the practical experiences and operational challenges of individuals directly affected by recent severe natural disasters in Vietnam. Observations from the unprecedented 2025 typhoon season, which recorded 21 storms and extreme rainfall peaks, indicate that the primary cause of communication failure during floods is the widespread disruption of local power grids. This disruption subsequently disables regional telecommunication transceivers, leading to intermittent or completely disconnected cellular networks. Victims frequently report that conventional messaging applications become unusable precisely when emergency communication is most critical. Consequently, stranded individuals are unable to transmit their exact coordinates, forcing rescue teams to conduct time-consuming, blind searches based on vague or outdated information. Furthermore, the rapid depletion of mobile device batteries during emergencies dictates a strict requirement: any proposed application must be capable of silently capturing and transmitting location data in the background with minimal energy consumption and bandwidth usage.

% Concurrently, an analysis of public sentiment regarding post-disaster relief operations highlights a severe crisis of trust in charity management. Following several high-profile controversies where tens of billions of VND in relief funds were mismanaged or lacked independent auditing, the community has articulated a strict demand for absolute financial transparency. Currently, donors encounter significant difficulties in verifying whether their financial contributions actually reach the intended recipients, as traditional methods of posting unstructured bank statements on social media are highly susceptible to manipulation and lack traceability. There is a clear, practical requirement from the community for a centralized, tamper-proof platform that standardizes the charity lifecycle. Donors and authorities require a system that not only publicly records financial transactions but also visually tracks the physical distribution of relief goods. This dual practical context—the need for resilient emergency communication and the demand for absolute charity transparency—directly dictates the core functional requirements of the proposed system.

\subsection{Analysis of existing systems and similar applications}
\label{subsection:2.1.1}

To identify the functional requirements of the proposed system, several existing applications related to communication, location sharing, emergency rescue, and charity management were studied. The selected systems include Zalo~\cite{zaloapp}, a web-based SOS platform~\cite{antoanmualu}, and Whoo~\cite{whooapp}, as they represent different approaches to location sharing and emergency support.

Zalo - instant messaging application: Zalo~\cite{zaloapp} is the most popular communication tool in Vietnam. Users can share their live location with contacts. However, this feature is strictly session-based and requires manual activation within a specific chat box. Crucially, it lacks continuous background tracking and entirely fails when the Internet connection drops. It also does not provide a centralized mechanism to broadcast SOS signals directly to authorized rescue teams.

"An toàn cứu hộ" - web-based SOS platform: It is dedicated platforms designed to transmit distress signals. While directly connected to rescue operations, its primary limitation is their reliance on web browsers, meaning they operate strictly online. If a victim loses network connectivity, the system immediately drops the connection, preventing any signal transmission. Furthermore, these platforms often lack features for victims to share their status with family members.

Whoo - social location app: Whoo~\cite{whooapp} is a mobile application designed for sharing locations among friends. It successfully implements a mechanism to save the user's last known location (last-location caching). However, it is primarily built for social networking rather than emergency response, meaning it cannot link victims with professional rescue units. Additionally, the application has been reported to suffer from synchronization instability.

Table \ref{table:2.1} summarizes the comparison between these existing systems and the proposed FloodHelper application. The comparison shows that existing systems typically focus on only one aspect of disaster support. Messaging applications mainly provide communication functions, while emergency platforms emphasize incident reporting. None of the surveyed systems integrates emergency rescue coordination with charity campaign management into a unified platform. Therefore, this graduation project proposes an integrated system that combines these functions to better support flood disaster response and post-disaster relief activities.

\begin{table}[htpb]
    \centering
    \caption{Comparison of existing applications and the proposed system}
    \label{table:2.1}
    \renewcommand{\arraystretch}{1.3} % Tăng khoảng cách các dòng để bảng nhìn thoáng hơn
    \begin{tabular}{|l|c|c|c|c|}
        \hline
        \textbf{Feature}               & \textbf{Zalo} & \textbf{Web SOS} & \textbf{Whoo} & \textbf{Proposed System} \\
        \hline
        Platform                       & Mobile / PC   & Web              & Mobile        & Mobile \& Web            \\
        \hline
        Offline Last-Location Saving   & No            & No               & Yes           & Yes                      \\
        \hline
        Continuous Background Tracking & No            & No               & Yes           & Yes                      \\
        \hline
        Direct Rescue Team Integration & No            & Yes              & No            & Yes                      \\
        \hline
        Transparent Charity Management & No            & No               & No            & Yes                      \\
        \hline
    \end{tabular}
\end{table}

\subsection{Functional requirements of the proposed system}
\label{subsection:2.1.2}

Based on the survey of existing applications and the practical requirements of flood disaster response, the proposed FloodHelper system is designed to provide an integrated platform that supports emergency communication, rescue coordination, and charity campaign management. The major software features to be developed are described as follows.

\subsubsection{User and role management}
\label{subsubsection:user_role_management}
The system provides user registration, authentication, and profile management for different types of users, including Normal Users, Benefactors, Rescuers, Authorities, and Administrators. Depending on the assigned role, users are granted different permissions to access the corresponding system functions. The system also supports role application and profile verification workflows to ensure that privileged roles are approved by the responsible authorities before becoming active.

\subsubsection{Friend connection and location sharing}
\label{subsubsection:friend_location_sharing}
The application enables users to establish friend connections and periodically share their current locations with connected friends while the application is running. Besides geographical coordinates, the shared information also reflects the user's current availability status, allowing friends to know whether the user is currently online or offline. To provide greater control over personal privacy, users can temporarily suspend location updates for selected friends, who will then only be able to view the user's last shared location until sharing is resumed.

\subsubsection{Emergency rescue support}
\label{subsubsection:emergency_rescue_support}
The system provides emergency communication functions that allow users to broadcast distress signals together with their current locations during flood situations. Rescue personnel can monitor active emergency requests, identify victims on a digital map, and update the handling status of each rescue request. Throughout the rescue process, location information is periodically synchronized so that both victims and rescuers can track rescue progress more effectively.

\subsubsection{Charity campaign management}
\label{subsubsection:charity_campaign_management}
The system supports the complete lifecycle of charity campaigns, from campaign creation and approval to donation collection, relief distribution, and financial reporting. Benefactors can organize charity campaigns, while authorities review submitted campaigns before publication. During campaign execution, donors are able to monitor campaign progress, distribution activities, and financial statements to improve transparency.

\subsubsection{Public information and notification}
\label{subsubsection:public_information_notification}
Authorities can publish public announcements regarding weather conditions, emergency instructions, evacuation guidance, or other disaster-related information. The application delivers these announcements to users located in the affected areas, allowing important information to be disseminated promptly. In addition, the system generates notifications for rescue activities, charity campaigns, and account-related events to keep users informed of important updates.

\subsubsection{Administrative management}
\label{subsubsection:administrative_management}
Administrative functions are provided to support overall system operation and supervision. Administrators can manage user accounts, while authorities review role applications, profile update requests, and charity campaign requests. These management functions help ensure that the platform operates in accordance with the defined policies and maintains the integrity of system data.

\section{Functional overview}
\label{section:2.2}
\subsection{General use case diagram}
\label{subsection:2.2.1}

Figure \ref{fig:Fig221} presents the overall use case diagram of the FloodHelper system. In details, the system involves five primary actors, each responsible for different operational tasks.

\begin{figure}[H]
    \centering
    \includegraphics[width=\linewidth]{Figure/UCDiagram/General Usecase.png}
    \caption{General usecase diagram}
    \label{fig:Fig221}
\end{figure}

Normal user is the primary user of the mobile application. A normal user can manage personal information, connect with friends, share locations, broadcast emergency distress signals, participate in charity campaigns, and receive public announcements.

Benefactor represents individuals or organizations responsible for organizing charity campaigns. In addition to the functions available to normal users, benefactors can create, manage, and update charity campaigns.

Rescuer is responsible for monitoring emergency requests and coordinating rescue activities. Rescuers receive distress signals, locate victims, and update rescue statuses during emergency operations.

Authority supervises rescue and charity activities. Authorities review campaign requests, approve user role applications, verify profile updates, and publish public announcements.

Administrator manages system operation by maintaining user accounts and administrative data to ensure stable platform operation.

\subsection{Decompose use case "manage my account"}
\label{subsection:2.2.2}

Figure \ref{fig:Fig222} presents the account management module of the FloodHelper system. The module provides common account functions, including registration, profile management, password management, and role application. While all users can manage their personal accounts, users with specialized roles are required to submit profile update requests for verification before changes are applied. This mechanism helps maintain the reliability of verified user information throughout the system.

\begin{figure}[H]
    \centering
    \includegraphics[width=\linewidth]{Figure/UCDiagram/Manage my account.png}
    \caption{Decompose use case "manage my account"}
    \label{fig:Fig222}
\end{figure}

\subsection{Decompose use case "interact with map"}
\label{subsection:2.2.3}

Figure \ref{fig:Fig223} illustrates the use cases related to map interaction, which are available to the Normal User.  This module provides functionalities for map navigation, location visualization, and privacy management.

\begin{figure}[H]
    \centering
    \includegraphics[width=\linewidth]{Figure/UCDiagram/Interact with map.png}
    \caption{Decompose use case "interact with map"}
    \label{fig:Fig223}
\end{figure}

Users can access both the "View transportation map" and the "View weather map" to obtain geographical and environmental information. In both cases, the corresponding map pins are displayed to present relevant locations and events on the map.

To support navigation, users can focus the map on their current location or on the location of a selected friend. The system also allows users to manage the visibility of shared information by freezing or unfreezing location sharing for specific friends. In addition, users may hide charity campaign map pins to reduce visual clutter and improve the readability of the map interface.

\subsection{Decompose use case "make friend"}
\label{subsection:2.2.4}

Figure \ref{fig:Fig224} presents the use cases related to friend management for the Normal User. This module enables users to establish and maintain friend relationships within the FloodHelper system.

\begin{figure}[H]
    \centering
    \includegraphics[width=\linewidth]{Figure/UCDiagram/Make Friend.png}
    \caption{Decompose use case "make friend"}
    \label{fig:Fig224}
\end{figure}

Users can search for other users and send friend requests to initiate a connection. The system also allows users to view pending and previously sent friend requests, as well as cancel requests that have not yet been processed. When receiving a friend request, users may either accept or reject the request. In addition, existing friendships can be removed when users no longer wish to share information with each other.

These functionalities establish the social network required for subsequent features such as location sharing, emergency communication, and rescue coordination.

\subsection{Decompose use case "manage charity campaign"}
\label{subsection:2.2.5}

Figure \ref{fig:Fig225} illustrates the use cases related to charity campaign management, involving three primary actors: Benefactor, Normal User, and Authority.

\begin{figure}[H]
    \centering
    \includegraphics[width=\linewidth]{Figure/UCDiagram/Manage charity campaign.png}
    \caption{Decompose use case "manage charity campaign"}
    \label{fig:Fig225}
\end{figure}

The Benefactor is responsible for organizing and managing charity campaigns. Benefactors can create, update, and delete campaign information, submit campaign requests for approval, publish campaign announcements, manage distribution items, and upload bank statements after the campaign is completed. These functions support the management of a charity campaign throughout its lifecycle.

The Normal User participates as a donor and can access campaign information through various viewing functions, including campaign details, announcements, distribution records, transaction history, and bank statements. Users can also donate to active campaigns and view campaign locations on the map to better understand the distribution of relief activities.

The Authority supervises charity campaigns by reviewing submitted campaign requests before publication. Authorities can approve or reject campaign requests and suspend active campaigns when necessary to ensure that charity activities comply with the platform's management policies.

\subsection{Decompose use case "manage distress signal"}
\label{subsection:2.2.6}

Figure \ref{fig:Fig226} presents the use cases related to emergency distress signal management, involving two primary actors: the Normal User and the Rescuer.

\begin{figure}[H]
    \centering
    \includegraphics[width=\linewidth]{Figure/UCDiagram/Manage distress signal.png}
    \caption{Decompose use case "manage distress signal"}
    \label{fig:Fig226}
\end{figure}

The Normal User can create, view, and update distress signals when emergency assistance is required. During signal creation, users provide essential information describing the emergency situation and their current location to facilitate rescue operations. If assistance is no longer needed, users may revoke an active distress signal before it is handled by a rescuer.

The Rescuer is responsible for monitoring and responding to active distress signals. Rescuers can view and sort broadcasting signals to identify and prioritize emergency requests, locate victims on the map, and handle selected distress signals throughout the rescue process. These functions support effective rescue coordination by enabling rescuers to monitor ongoing incidents and respond to emergency requests in a timely manner.

\subsection{Decompose use case "manage user"}
\label{subsection:2.2.7}

Figure \ref{fig:Fig227} presents the use cases related to user management, which are performed exclusively by the Admin. This module provides administrative functions for managing user accounts and maintaining system operation.

\begin{figure}[H]
    \centering
    \includegraphics[width=\linewidth]{Figure/UCDiagram/Manage user.png}
    \caption{Decompose use case "manage user"}
    \label{fig:Fig227}
\end{figure}

The Admin can search for users by email to quickly locate specific accounts within the system. To enforce platform policies and manage inappropriate user activities, the Admin is authorized to ban or unban user accounts when necessary. In addition, the Admin is responsible for managing Authority accounts by creating, viewing, and updating their information. These functions support the administration of user access and ensure that authorized personnel are properly maintained within the system.

\subsection{Decompose use case "review role requests"}
\label{subsection:2.2.8}

Figure \ref{fig:Fig228} presents the use cases related to role request management, which are performed by the Authority. This module enables the Authority to review requests submitted by users who wish to obtain specialized roles within the FloodHelper system.

\begin{figure}[H]
    \centering
    \includegraphics[width=\linewidth]{Figure/UCDiagram/Review role requests.png}
    \caption{Decompose use case "review role requests"}
    \label{fig:Fig228}
\end{figure}

The Authority can view role requests based on their current processing status to facilitate the review process. After evaluating each request, the Authority may either approve or reject the pending role request. An approved request grants the applicant the corresponding system role, while a rejected request prevents the requested role from being assigned.

These functions ensure that specialized roles, such as Benefactor and Rescuer, are assigned only after verification by the responsible authority.

\subsection{Decompose use case "review updating-profile requests"}
\label{subsection:2.2.9}

Figure \ref{fig:Fig229} presents the use cases related to profile-updating request management, which are performed by the Authority. This module enables the Authority to review profile modification requests submitted by users with specialized roles, such as Benefactors and Rescuers.

\begin{figure}[H]
    \centering
    \includegraphics[width=\linewidth]{Figure/UCDiagram/Review updating-profile requests.png}
    \caption{Decompose use case "review updating-profile requests"}
    \label{fig:Fig229}
\end{figure}

The Authority can view profile-updating requests based on their current processing status to facilitate the review process. After examining the submitted information, the Authority may either approve or reject each pending request. Approved requests update the corresponding user profile with the verified information, while rejected requests leave the existing profile unchanged.

These functions help ensure that the information displayed for specialized users remains accurate and verified before being made available within the system.

\subsection{Business process}
\label{subsection:2.2.10}

The FloodHelper system includes several core business processes that involve interactions among multiple actors to accomplish major system functionalities. These processes describe the sequence of activities performed by users and the system in typical operational scenarios, providing a high-level view of how different functional modules collaborate to achieve the intended objectives.

To illustrate these workflows, business process diagrams are used to represent the interactions between actors and the system. In this project, three representative business processes are presented: Update Personal Information, Broadcast Signal and Rescue Victim, and Create and Manage Charity Campaign. These processes cover the key operational scenarios of user profile management, emergency rescue coordination, and charity campaign management within the FloodHelper system.

The following subsections describe each business process in detail.

\subsubsection{Update personal information}

Figure \ref{fig:Fig2210a} illustrates the business process for updating personal information. The process begins when a user submits a request to modify their profile information. Depending on the user's role, the submitted information may be updated directly or require verification by the Authority before the changes are applied. This process ensures that profile information remains accurate while maintaining the reliability of verified user accounts.

\begin{center}
    \includegraphics[width=1.0\linewidth]{Figure/BusinessProcess/BusinessProcess_UpdatePersonalInformation.png}
    \captionof{figure}{Business process "Update personal information"}
    \label{fig:Fig2210a}
\end{center}

\subsubsection{Broadcast signal and rescue victim}

Figure \ref{fig:Fig2210b} presents the business process for emergency rescue coordination. The workflow starts when a user broadcasts a distress signal containing the required emergency information and current location. The system makes the distress signal available to rescuers, who can monitor active requests, locate the victim, and handle the assigned rescue task. The process concludes when the rescue request is completed or the distress signal is revoked.

\begin{center}
    \includegraphics[width=1.0\linewidth]{Figure/BusinessProcess/BusinessProcess_BroadcastSignalAndRescueVictim.png}
    \captionof{figure}{Bussiness process "Broadcast signal and rescue victim"}
    \label{fig:Fig2210b}
\end{center}

\subsubsection{Create and manage charity campaign}

Figure \ref{fig:Fig2210c} illustrates the business process for charity campaign management. The workflow begins with a Benefactor creating a charity campaign and submitting it for approval. After the campaign is approved by the Authority, it becomes available for public donations. Throughout the campaign, the Benefactor manages announcements, distribution activities, and financial reports, while donors can contribute to the campaign and monitor its progress until the campaign is completed.

\begin{center}
    \includegraphics[width=1.0\linewidth]{Figure/BusinessProcess/BusinessProcess_CreateAndManageCharityCampaign.png}
    \captionof{figure}{Business process "Create and manage charity campaign"}
    \label{fig:Fig2210c}
\end{center}

\section{Functional description}
\label{section:2.3}

\subsection{Description of use case "Broadcast signal"}
\label{section:2.3.1}

Table \ref{table:uc_broadcast_signal} presents the detailed specification of the use case "broadcast signal".

\renewcommand{\arraystretch}{1.3}
\begin{xltabular}{\textwidth}{|p{3.5cm}|X|}
    \caption{Use case specification for "Broadcast signal"}
    \label{table:uc_broadcast_signal} \\
    \hline
    \endfirsthead

    \multicolumn{2}{c}{{\tablename\ \thetable\ -- Continued from previous page}} \\
    \hline
    \endhead

    \hline
    \multicolumn{2}{r}{\textit{Continued on next page}} \\
    \endfoot

    \hline
    \endlastfoot

    Use case name & Broadcast signal \\
    \hline
    Use case ID & UC001-BS \\
    \hline
    Actors & Normal user, Rescuer \\
    \hline
    Pre-conditions &
    - The user has successfully logged into the system. \newline
    - The user's location service is available. \newline
    - The user has not created another active distress signal. \\
    \hline
    Main flow &
    1. The user selects the Broadcast Distress Signal function. \newline
    2. The system displays the distress signal form. \newline
    3. The user enters the required emergency information. \newline
    4. The user submits the distress signal. \newline
    5. The system validates the submitted information. \newline
    6. The system creates a new distress signal and records its information. \\
    \hline
    Alternative and exception flows &
    -Alternative flow 3.1: At Step 3, the user does not enter any emergency information. The system applies default values for the distress signal and continues with Step 4. \newline
    -Alternative flow 4.1: At any time before Step 4, the user cancels the operation. The system discards the entered information and returns to the previous screen. \\
    \hline
    Post-conditions &
    - A new distress signal is successfully created. \newline
    - The distress signal is available for rescuers to monitor. \newline
    - The user can track the status of the submitted distress signal. \\
    \hline
\end{xltabular}

\subsection{Description of use case "Handle signal"}
\label{section:2.3.2}

Table \ref{table:uc_handle_signal} presents the detailed specification of the use case "handle signal".

\renewcommand{\arraystretch}{1.3}
\begin{xltabular}{\textwidth}{|p{3.5cm}|X|}
    \caption{Use case specification for "Handle signal"}
    \label{table:uc_handle_signal} \\
    \hline
    \endfirsthead

    \multicolumn{2}{c}{{\tablename\ \thetable\ -- Continued from previous page}} \\
    \hline
    \endhead

    \hline
    \multicolumn{2}{r}{\textit{Continued on next page}} \\
    \endfoot

    \hline
    \endlastfoot

    Use case name & Handle signal \\
    \hline
    Use case ID & UC002-HS \\
    \hline
    Actors & Rescuer \\
    \hline
    Pre-conditions &
    - The Rescuer has successfully logged into the system. \newline
    - The Rescuer is viewing the map with active distress signals. \newline
    - A victim has broadcasted an active distress signal. \\
    \hline
    Main flow &
    1. The Rescuer selects an active distress signal on the map. \newline
    2. The system displays the details of the selected distress signal. \newline
    3. The Rescuer confirms handling the distress signal. \newline
    4. The system validates the request and updates the status of the distress signal to "Being handled". \newline
    5. The system notifies the victim that their distress signal is being handled. \newline
    6. The system updates the map view for other rescuers to show that the signal is being handled. \\
    \hline
    Alternative and exception flows &
    -Alternative flow 3.1: At Step 3, the Rescuer decides not to handle the signal and cancels the operation. The system returns to the map view. \newline
    -Alternative flow 5.1: At Step 5, if the victim is also registered as a rescuer, the system updates their personal signal status while maintaining their access to the rescue coordination map. \\
    \hline
    Post-conditions &
    - The distress signal status is successfully updated to "Being handled". \newline
    - The victim is notified of the rescue progress. \newline
    - Other rescuers are updated to avoid duplicate rescue deployments. \\
    \hline
\end{xltabular}

\subsection{Description of use case "Revoke signal"}
\label{section:2.3.3}

Table \ref{table:uc_revoke_signal} presents the detailed specification of the use case "revoke signal".

\renewcommand{\arraystretch}{1.3}
\begin{xltabular}{\textwidth}{|p{3.5cm}|X|}
    \caption{Use case specification for "Revoke signal"}
    \label{table:uc_revoke_signal} \\
    \hline
    \endfirsthead

    \multicolumn{2}{c}{{\tablename\ \thetable\ -- Continued from previous page}} \\
    \hline
    \endhead

    \hline
    \multicolumn{2}{r}{\textit{Continued on next page}} \\
    \endfoot

    \hline
    \endlastfoot

    Use case name & Revoke signal \\
    \hline
    Use case ID & UC003-RS \\
    \hline
    Actors & Normal user, Rescuer \\
    \hline
    Pre-conditions &
    - The user has successfully logged into the system. \newline
    - The user has an active distress signal. \\
    \hline
    Main flow &
    1. The user selects the option to revoke their active distress signal. \newline
    2. The system requests confirmation from the user. \newline
    3. The user confirms the revocation. \newline
    4. The system validates the request and marks the distress signal as inactive. \newline
    5. The system stops broadcasting the distress signal. \newline
    6. The system notifies rescuers in the area that the distress signal is no longer active. \\
    \hline
    Alternative and exception flows &
    -Alternative flow 3.1: At Step 3, the user decides to cancel the revocation. The system returns the user to the active signal details view. \\
    \hline
    Post-conditions &
    - The distress signal status is successfully marked as inactive. \newline
    - Rescuers in the area are updated, and the distress signal is cleared from their coordination maps. \\
    \hline
\end{xltabular}



\subsection{Description of use case "Publish public announcement"}
\label{section:2.3.4}

Table \ref{table:uc_publish_announcement} presents the detailed specification of the use case "publish public announcement".

\renewcommand{\arraystretch}{1.3}
\begin{xltabular}{\textwidth}{|p{3.5cm}|X|}
    \caption{Use case specification for "Publish public announcement"}
    \label{table:uc_publish_announcement} \\
    \hline
    \endfirsthead

    \multicolumn{2}{c}{{\tablename\ \thetable\ -- Continued from previous page}} \\
    \hline
    \endhead

    \hline
    \multicolumn{2}{r}{\textit{Continued on next page}} \\
    \endfoot

    \hline
    \endlastfoot

    Use case name & Publish public announcement \\
    \hline
    Use case ID & UC004-PAP \\
    \hline
    Actors & Authority \\
    \hline
    Pre-conditions &
    - The Authority is logged into the system. \newline
    - The Authority has the permissions to publish announcements. \\
    \hline
    Main flow &
    1. The Authority selects the option to create a public announcement. \newline
    2. The system displays the announcement creation form. \newline
    3. The Authority enters the announcement details (such as title, content, and optional file attachments) and submits the form. \newline
    4. The system validates the input. \newline
    5. The system saves the public announcement. \newline
    6. The system notifies all users in the Authority's ward of the new announcement. \newline
    7. The system displays a confirmation message to the Authority and refreshes the announcements list. \\
    \hline
    Alternative and exception flows &
    -Alternative flow 3.1: At Step 3, the Authority decides to cancel the creation. The system discards the input and returns to the previous screen. \newline
    -Alternative flow 3.2: At Step 3, if the Authority attaches files, the system uploads and stores the files associated with the announcement. \\
    \hline
    Post-conditions &
    - The public announcement is successfully published and stored. \newline
    - Users in the designated ward are notified. \newline
    - The announcement list is updated for the Authority and users. \\
    \hline
\end{xltabular}

\section{Non-functional requirement}
\label{section:2.4}
Performance (Latency): Real-time location tracking and signal broadcasting via the MQTT broker must maintain a delivery latency of under 500ms to ensure accurate, up-to-date mapping of victims and rescuers.

Reliability (Network Resilience): The mobile application must utilize delay-tolerant networking strategies and background isolates to queue emergency signals and location data locally, ensuring eventual transmission even when operating in unstable or heavily degraded network environments.

Scalability (Concurrency): The containerized backend infrastructure must be capable of automatically scaling to handle sudden, massive spikes in concurrent connections and API requests during widespread emergency events without system degradation.

Maintainability (Architecture): The system's codebase must strictly adhere to clean architecture principles (such as the Model-View-ViewModel (MVVM) pattern for the frontend) to ensure straightforward unit testing, predictable state management, and the seamless integration of future communication protocols.

%%%%%%%%%%%%%%%%%%%%%%%%%%%%%%%%%%%

\end{document}